
\definecolor{bdcolor}{RGB}{68,170,168} % define border color
\newcommand{\overviewBox}[1]{%
\tikz[baseline=-0.5ex]{\node[rounded corners=0.2em, 
fill=white, 
draw=bdcolor, 
minimum size=0.9em, 
text=bdcolor, 
text centered,
inner sep=0pt,
line width=0.9pt,
font=\fontsize{9pt}{0.9em}\fontseries{ul}\selectfont](TsNode){#1}}}

\section{Introduction}

Large Language Models (LLMs) have rapidly gained prominence in recent years, demonstrating significant value and potential across various sectors such as industry\cite{tiro2023possibility,wang2023chatgpt}, healthcare\cite{thirunavukarasu2023large,abd2023large,sallam2023chatgpt}, science\cite{wardat2023chatgpt,agler2011waste}, and Finance\cite{wu2023bloomberggpt,yang2023fingpt}. Despite their unprecedented language comprehension and text generation capabilities, LLMs can face limitations when dealing with domain-specific knowledge, real-time updated information, and proprietary data, which are not in their pre-training corpus\cite{gao2023retrieval}. A promising paradigm for addressing this limitation is Retrieval-Augmented Generation (RAG), which integrates a retrieval component within the generation process to leverage the information contained in a customized large text corpus. This process not only enriches the contextual depth of the responses but also boosts their factual accuracy and specificity \cite{peng2024graph, fan2024survey}. Although RAG has achieved impressive results, it faces limitations in real-world scenarios that involve complex structures of relationships among different entities in the customized text corpus \cite{peng2024graph, procko2024graph}. To address this, recent works \cite{edge2024local,wu2024medical,guo2024lightrag} propose introducing graphs as an intermediate representation during the retrieval process (i.e. Graph-based RAG) to capture better and leverage the structured relational knowledge contained in the customized text corpus, bringing about new opportunities to elevate the precision and comprehensiveness of generation results in RAG applications.

However, despite the promising prospects of the Graph-based RAG (GraphRAG) paradigm, current RAG application developers often face challenges in analyzing the effectiveness of GraphRAG on their datasets. This reduces their confidence and motivation to deploy GraphRAG solutions in actual production. This difficulty arises from the two additional layers of complexity introduced by GraphRAG compared to traditional RAG: First, the information processing pipeline of a GraphRAG system is much longer and more complex than that of traditional RAG. Starting from raw text, the information goes through a series of intricate processes, including entity and relation extraction, graph construction, information aggregation and summarization, and parallel multi-stage retrievals. Users lack effective tools to track and analyze these processes. Second, the information processing in a GraphRAG system involves frequent and extensive calls to LLMs. It is difficult to understand the context and role of each LLM invocation within the graph during the analysis process.


As such, it is highly desirable to provide analysis support for GraphRAG developers to examine the GraphRAG process on their customized text corpus, enabling them to make informed design decisions and adjustments. To clarify the analysis targets, we survey recent GraphRAG-related papers and open-source projects to abstract a common pipeline for GraphRAG. For the sake of better understanding the underlying challenges in the analysis process, we conducted a formative study with 5 experienced RAG experts to have a grasp of the current analysis process, common analysis strategies, and the challenges they encounter when analyzing the GraphRAG process.

Based on the findings of the formative study, we propose a two-stage visual analysis framework for the process diagnosis of GraphRAG systems. In the first stage, it leverages the power of the LLM to provide automatic evaluation of model answers and identify the suspicious incorrect recalls that mislead the inference process. In the second stage, it allows users to trace the causes of the incorrect recalls through multi-facet relevance analysis on the graph. Based on this framework, we design \textit{XGraphRAG}, a visual analysis system consisting of four interrelated views. The QA View (Fig~\ref{fig:teaser}A) and the Inference Trace View (Fig~\ref{fig:teaser}B) serve as the springboard for exploration that helps users compare the discrepancies between the answer and ground truth to identify suspicious recalls. The Topic Explore View (Fig~\ref{fig:teaser}C) and the Entity Explore view (Fig~\ref{fig:teaser}D) allow users to conduct relevance analysis on the graph from different aspects. The LLM Invocation View (Fig~\ref{fig:teaser}E) adaptively displays structured details of the LLM invocation related to the intermediate output the user is interested in during the analysis process. We conduct a usage scenario demonstration and a user study to verify the effectiveness of our approach.

In summary, this paper makes the following contributions:
\begin{itemize}
\item We identify challenges in the analysis of the GraphRAG process and distill design requirements.
\item We propose a visual analysis framework that allows GraphRAG developers to examine the GraphRAG process efficiently and systematically.
\item We present \textit{XGraphRAG}, a prototype system that instantiates our framework, whose effectiveness is verified with a usage scenario demonstration and a user study.
\end{itemize}